\chapter{基板の開発}
\section{概要}
衛星システムは複数の電子回路基板によって構成され、それらが電気的に結合することで機能する。本開発は、機能実証を行うための試作機であるエンジニアリングモデル（EM）の開発と、実機搭載用のフライトモデル（FM）の開発の大きく分けて二段階を経て実施された。以下に、開発した基板とその開発プロセスの概要を述べる。

\subsection{基板の概要}
磁気トルカボード、（コネクタボード）、姿勢系システム搭載ボード、赤外線センサインターフェースボード、Spresense拡張ボード

\subsection{基板設計ソフトウェア EAGLE}
本研究の基板開発では、Autodesk社が提供するプリント基板設計者向けの電子設計オートメーションソフトウェア「EAGLE」を用いた。EAGLEで基板設計を行った後、そのガーバーデータを基板製造会社であるP板.comに提出して基板製造を依頼することで、自身の設計した基板が完成する。

Autodesk社のEAGLEは2026年6月7日をもって販売およびサポートを終了することが決定しており、その後は同社Autodesk Fusion内のAutodesk Fusion Electronicsで同様の基板設計ができる。本研究で開発した主要な基板のEAGLEデータについては、Autodesk Fusion Electronicsに移行して〇〇に置いてある。% memo 保存場所を書く

\subsection{開発プロセスの概要}
FM基板の開発には、EM基板の課題を受けた設計変更と要求
基板の開発は、設計要求の確認、基板仕様の検討、回路設計、アートワーク設計、基板の評価という順で行う。以下にそれぞれの詳細をまとめる。

\subsubsection{設計要求の確認}
\subsubsection{基板仕様の検討}
\subsubsection{回路設計}
\subsubsection{アートワーク設計}
\subsubsection{基板の評価}

\section{エンジニアリングモデルの基板開発}
\subsection{基板仕様}
\subsection{課題とその究明}


\section{フライトモデルの基板開発}
\subsection{EMの課題を受けて}
\subsection{設計要求}
\subsubsection{磁気トルカボード}
\subsubsection{磁気トルカコネクタボード}
\subsubsection{姿勢系システム搭載ボード}
\subsubsection{赤外線センサインターフェースボード}
\subsubsection{Spresense拡張ボード}


\subsection{基板仕様の検討}
\subsubsection{磁気トルカボード}
\subsubsection{磁気トルカコネクタボード}
\subsubsection{姿勢系システム搭載ボード}
\subsubsection{赤外線センサインターフェースボード}
\subsubsection{Spresense拡張ボード}

\subsection{回路設計}
\subsubsection{磁気トルカボード}
\subsubsection{磁気トルカコネクタボード}
\subsubsection{姿勢系システム搭載ボード}
\subsubsection{赤外線センサインターフェースボード}
\subsubsection{Spresense拡張ボード}

\subsection{アートワーク設計}
\subsubsection{磁気トルカボード}
\subsubsection{磁気トルカコネクタボード}
\subsubsection{姿勢系システム搭載ボード}
\subsubsection{赤外線センサインターフェースボード}
\subsubsection{Spresense拡張ボード}

\subsection{基板の評価}
\subsubsection{磁気トルカボード}
\subsubsection{磁気トルカコネクタボード}
\subsubsection{姿勢系システム搭載ボード}
\subsubsection{赤外線センサインターフェースボード}
\subsubsection{Spresense拡張ボード}

\section{（仮称）ハードウェア設計の注意点}



% memo ハーネスのことも述べる